
\begin{titlepage}

	\begin{center}
	% Partie supérieure de la page
	\includegraphics[width=0.3\textwidth]{./Images/Logos/fs.pdf}\hfill
	\includegraphics[width=0.5\textwidth]{./Images/Logos/umons_charleroi.png}\\[1cm]
	
	{\LARGE Université de Mons}\\[0.5cm]
	{\Large Faculté des Sciences}\\[1cm]
	
	{\Large US-MC-INFO60-014-C — Lecture et rédaction scientifique}\\[1cm]
	
	% Titre
	\HRule \\[0.4cm]
	{ \Huge \bfseries TOR:}\\[0.2cm]
	{ \LARGE \bfseries Le Routage Onion de Seconde Génération}\\[0.4cm]
	\HRule \\[1.5cm]
	
	% Auteur et directeur
	\begin{minipage}{0.4\textwidth}
	\begin{flushleft} \large
	\emph{Auteur:}\\
	Baptiste \textsc{Grosjean}
	\end{flushleft}
	\end{minipage}
	~
	\begin{minipage}{0.4\textwidth}
	\begin{flushright} \large
	\emph{Directeur:}\\
	Alain \textsc{Buys}
	\end{flushright}
	\end{minipage}\\[2cm]
	
	\begin{minipage}{0.4\textwidth}
	\begin{flushleft} \large
	\emph{Numéro d'étudiant:}\\
	232732
	\end{flushleft}
	\end{minipage}\\[2cm]
	
	% Résumé
	\renewcommand{\abstractname}{Résumé}
	\begin{abstract}
	\acrfull{tor} est un système de communication sécurisé et fiable. Ce travail de rédaction scientifique vise à fournir une compréhension détaillée du protocole \acrshort{tor}, de son fonctionnement et de son importance dans les réseaux modernes. Nous utiliserons \LaTeX\ pour produire une documentation claire et précise du protocole ainsi que de son implémentation pratique.
	\end{abstract}
	
	\vfill
	
	% Bas de la page
	{\large Juin 2024}
	\end{center}
	\end{titlepage}